\subsection{Conclusiones generales con evidencia cuantitativa}
\begin{enumerate}
    \item Los resultados del benchmark respaldaron la hipótesis inicial de que nuestra implementación con sockets TCP presentaría una menor latencia media que gRPC en las condiciones evaluadas. Con mensajes de 100 bytes, TCP registró 0,122\,ms frente a 0,547\,ms de gRPC, mientras que con cargas de 10\,KiB los valores fueron de 0,197\,ms y 0,564\,ms, respectivamente. Esta diferencia también se observó en las medianas y los percentiles 95. Sin embargo, los resultados corresponden a las implementaciones del equipo y no demuestran que TCP sea más rápido en cualquier entorno distribuido.

    \item En cuanto al rendimiento, TCP también presentó valores superiores en el experimento secuencial. Con cargas de 100 bytes alcanzó una media de 576,43 mensajes por segundo, frente a 442,64 en gRPC. Para 10\,KiB, se registraron 533,24 y 431,98 mensajes por segundo, respectivamente. Estas cifras representan un rendimiento aproximadamente 30,2\,\% y 23,4\,\% mayor para TCP en cada tamaño evaluado. No obstante, la elección de un mecanismo de comunicación también debe considerar el tipado de datos, el mantenimiento y el manejo de errores, aspectos que no quedan reflejados únicamente en estas mediciones.

    \item La integración de los relojes de Lamport permitió registrar y verificar el orden causal de los eventos durante los intercambios entre tres nodos. Las corridas funcionales de TCP y gRPC completaron veinte solicitudes por mecanismo, con 80 eventos registrados y ningún error causal reportado. Además, en las cuatro pruebas concurrentes se registraron 300 solicitudes y 1200 eventos por condición, también sin errores causales detectados. Esto respalda el funcionamiento de las reglas del reloj lógico sobre los registros evaluados, aunque no demuestra tolerancia ante todos los posibles fallos de comunicación.

    \item Finalmente, se comprobó la importancia de definir condiciones experimentales claras para interpretar correctamente los resultados. El benchmark se realizó sobre localhost, con mensajes secuenciales y dos tamaños de carga útil, mientras que las pruebas concurrentes se desarrollaron por separado en otro entorno. Por esta razón, sus resultados no deben mezclarse para comparar latencias o rendimiento. Como mejora futura, sería conveniente repetir las mediciones en una red física, evaluar diferentes niveles de concurrencia y aplicar pruebas controladas de caída de nodos. Esto permitiría ampliar el análisis de rendimiento y fiabilidad.

\end{enumerate}

\subsection{Conclusiones individuales}
\pendiente{Según la guía: al menos 150 palabras POR INTEGRANTE, con contenido propio y relacionado con su contribución; no copiarlas entre sí.}
\subsubsection{Joseph}
% Conclusión individual: Joseph
En mi participación me enfoqué en la comunicación mediante sockets TCP y en cómo estructurar el intercambio entre el cliente y el servidor. La implementación usa una conexión que se reutiliza para varios mensajes; cada solicitud viaja como JSON codificado en UTF-8 y termina con un salto de línea para que el servidor pueda separar los mensajes. Al trabajar con este mecanismo entendí que TCP entrega un flujo de bytes y que, por eso, el formato, la delimitación y varias validaciones deben resolverse en la aplicación.

Las pruebas funcionales permitieron comprobar el intercambio de solicitudes y acuses entre los nodos. También aprendí que una conexión puede fallar o agotarse el tiempo de espera, y que esos casos deben distinguirse de una respuesta exitosa. En el benchmark secuencial de localhost, TCP registró menor latencia media que gRPC para los dos tamaños evaluados; aun así, ese resultado corresponde a las implementaciones y condiciones usadas en esta práctica, no a una regla para cualquier sistema.

La comparación me ayudó a entender que elegir un mecanismo de comunicación no depende solo de la velocidad. TCP ofreció control sobre el formato, pero exigió más trabajo manual para validar mensajes y manejar errores. Esta experiencia me permitió relacionar las decisiones de implementación con los resultados medidos y reconocer qué aspectos todavía habría que probar en una red y con otras cargas.
\subsubsection{Mera}
% Conclusión individual: Mera
Mi parte de la práctica fue la implementación con gRPC: escribí el contrato \texttt{messaging.proto}, el servidor y el cliente, y los integré con el reloj de Lamport de Roger. Una decisión bastanete importante fue que mi cliente tuviera la misma interfaz que el de TCP (\texttt{conectar}, \texttt{enviar}, \texttt{cerrar}), para que el compañero Mario pudiera medir los dos mecanismos con el mismo script. También decidí abrir un solo canal y esperar a que estuviera listo antes de medir, para no contar la conexión dentro de la latencia.

El problema que más me costó fue que el servidor seguía respondiendo con 6 en vez de 7 después de integrar el reloj; el código estaba bien, pero el archivo no se había guardado. En los 20 intercambios obtuve la misma secuencia de relojes que en TCP (acuse final 79 y reloj local 80), lo que me confirmó que el transporte no cambia el orden de Lamport. De esa manera aprendí que el tipado de gRPC facilita el trabajo, pero no reemplaza la validación: una solicitud sin \texttt{timestamp} llega con valor 0 sin dar error. Con 100 bytes, gRPC fue unas 4,5 veces más lento que TCP, así que en el PFC lo usaría donde importe más un contrato claro entre servicios que la latencia mínima.
\subsubsection{Roger}
% Borrador aportado por Roger; precisión y formato revisados con asistencia de IA.
Durante esta práctica, mi responsabilidad fue trabajar con el reloj lógico de
Lamport, integrar su funcionamiento en cada nodo, registrar los eventos y verificar
los resultados. Aprendí que cada nodo necesita su propio reloj porque los procesos
funcionan de manera independiente. También comprendí que el servidor y el cliente
de un mismo nodo deben compartir el contador para mantener un orden coherente.
Antes no tenía tan claro cómo se relacionaban los eventos entre distintos procesos,
pero ahora entiendo mejor el concepto de causalidad.

Una de las dificultades fue diferenciar un error del programa de un fallo esperado
durante las pruebas. Cuando se detuvo el nodo C, observé que B registró el envío
fallido y que su reloj avanzó de 7 a 8. Al revisar los registros, comprobé que la
comunicación entre A y B continuaba funcionando, lo que me permitió entender mejor
el comportamiento del sistema.

En cada corrida funcional de TCP y gRPC se completaron 20 solicitudes y 20 acuses,
con 80 eventos y 40 pares, sin errores de causalidad detectados por el verificador.
Esta experiencia me permitió comprender que Lamport ordena eventos, pero no
sincroniza la hora real. Considero que una mejora sería automatizar más
verificaciones para facilitar la detección de errores.

\subsubsection{Mario}
% Conclusión individual: Mario
\pendiente{Redactar mínimo 150 palabras de autoría propia. Referir decisiones, problemas concretos, pruebas y aprendizajes comprobables de la participación individual.}

